\section{Introduction}

\IEEEPARstart{C}{ardiovascular} disease accounts for approximately 17.9
million deaths annually~\cite{who2023cvd}, and automated electrocardiogram
(ECG) interpretation has advanced to cardiologist-level agreement on
standardised corpora~\cite{hannun2019,ribeiro2020}. Translating those models
into continuous monitoring, however, runs into a mundane engineering
constraint: the devices that would carry them differ by orders of magnitude in
memory, bandwidth and power. A smartwatch, a phone and a hospital workstation
cannot host the same artefact, and the prevailing response is to train, tune,
validate and separately certify one model per target~\cite{rajpurkar2022}.

Existing compression techniques inherit that structure. Knowledge
distillation~\cite{hinton2015}, pruning~\cite{han2016} and
quantisation~\cite{jacob2018} each yield a \emph{new} model for each operating
point; slimmable networks~\cite{yu2019} and early-exit
architectures~\cite{teerapittayanon2016} relax it partially but require
purpose-built multi-branch designs. Matryoshka Representation Learning
(MRL)~\cite{kusupati2022}, a concrete instance of the nested learning paradigm
recently formalised by Behrouz \etal~\cite{behrouz2025}, takes a different
route: it trains a single encoder such that every prefix $z_{1:m}$ of the
embedding is independently usable. One set of weights, one validation
pathway, and an operating point selected at deployment time by a pointer
offset. MRL underpins production embedding APIs~\cite{openai2024embeddings,
google2025gemma3n} but has not been examined for physiological time series.

\subsection{Scope, and what this paper does not claim}

An earlier version of this work framed the contribution as spanning a
``wearable-to-cloud continuum''. We have narrowed that framing deliberately,
because the evidence does not reach it and the distinction matters for
readers deciding whether to build on the method.

PTB-XL comprises \emph{clinical, twelve-lead, resting} ECGs recorded with wet
electrodes under supervision. Wearable acquisition differs along nearly every
axis that matters: one or two dry-electrode leads, ambulatory motion
artefacts, intermittent contact, and different noise statistics. Results
obtained on the former do not automatically transfer to the latter, and
presenting a device-tier table as though they do overstates what has been
shown. We therefore separate three claims that are easily conflated:

\begin{enumerate}
\item \textbf{Embedding compressibility} (demonstrated here). Diagnostic
information in twelve-lead ECG representations concentrates into a small
leading prefix, and nesting exploits this at negligible cost.
\item \textbf{Deployment mapping} (a design implication, not a result). A
given embedding budget admits a given truncation point; this is arithmetic,
and we present it as such.
\item \textbf{On-device wearable validation} (\emph{not} established). This
requires wearable-acquired data and microcontroller execution. We report
reduced-lead and artefact-corrupted proxies, which bound the question without
settling it, and identify it as the necessary next step.
\end{enumerate}

We also decline the conventional framing of accuracy superiority. Our models
do not exceed the strongest published PTB-XL classifiers, and the mechanism of
interest---nesting---is orthogonal to backbone quality. The contribution is
that adaptivity is obtainable at negligible cost, which is a statement about
the \emph{shape} of the accuracy--dimension curve rather than its height.

\subsection{Contributions}

\begin{enumerate}
\item \textbf{Nested representation learning for biosignals.} The first
application of Matryoshka representations to ECG classification, evaluated
across two backbone families, six nesting granularities and five training
seeds.

\item \textbf{Like-for-like baselines.} Fixed-dimension models trained with
the \emph{same} backbone, schedule, augmentation and seeds at every nesting
dimension---the comparison that isolates the cost of nesting, and which
cross-architecture comparisons cannot.

\item \textbf{Statistical rather than visual conclusions.} Dimension
invariance is assessed by an equivalence test with a pre-registered margin,
model comparisons by paired bootstrap and DeLong tests with Holm correction.
A claim of ``no difference'' requires an equivalence procedure; superiority
tests cannot supply it.

\item \textbf{Measured computation.} Latency, energy and memory are measured
on real hardware across GPU, CPU, quantised CPU and ONNX Runtime. The
resulting picture---latency flat in $d$---is less flattering than an
extrapolated one but is what practitioners need in order to budget correctly.

\item \textbf{Probing the hierarchy.} Ridge probes recovering eleven
physiological descriptors from each prefix convert the intuitive
coarse-to-fine story into a falsifiable measurement, and we report the outcome
whichever way it falls.

\item \textbf{Generalisation stress tests.} Zero-shot transfer to three
external corpora, reduced-lead ablation down to single-lead input, and
robustness to baseline wander, EMG, electrode motion and mains interference at
controlled SNR.
\end{enumerate}
